% !TeX program = xelatex
\documentclass[tikz, border=20pt]{standalone}

% 中文排版支持
\usepackage[UTF8]{ctex}

% 顶级期刊标准字体 (Times Roman 风格)
\usepackage{newtxtext, newtxmath}

% 引入必要的 TikZ 库
\usetikzlibrary{shapes.geometric, arrows.meta, positioning, calc, fit, backgrounds, shadows.blur}

\begin{document}

\begin{tikzpicture}[
    node distance=1.5cm and 2cm,
    font=\small, 
    >=Stealth,
    
    % 核心组件样式 (稍微加高以容纳举例)
    process/.style={rectangle, minimum width=4.5cm, minimum height=1.2cm, text centered, draw=black!80, fill=blue!8, rounded corners=4pt, blur shadow={shadow blur steps=5, shadow opacity=20}},
    decision/.style={diamond, aspect=2.2, minimum width=4cm, minimum height=1cm, text centered, draw=black!80, fill=yellow!12, blur shadow={shadow blur steps=5, shadow opacity=20}},
    
    % 奖惩节点：改为两端对齐/居左模式，方便分行展示说明与得分
    penalty/.style={rectangle, text width=3.8cm, minimum height=1.4cm, draw=red!60, fill=red!10, rounded corners=4pt, font=\small, blur shadow={shadow blur steps=5, shadow opacity=20}, align=left, inner sep=8pt},
    reward/.style={rectangle, text width=3.8cm, minimum height=1.4cm, draw=green!60!black, fill=green!10, rounded corners=4pt, font=\small, blur shadow={shadow blur steps=5, shadow opacity=20}, align=left, inner sep=8pt},
    dense/.style={rectangle, text width=4.0cm, minimum height=1.4cm, draw=orange!80, fill=orange!10, rounded corners=4pt, font=\small, blur shadow={shadow blur steps=5, shadow opacity=20}, align=left, inner sep=8pt},
    
    % 连接线样式
    arrow/.style={thick, ->, draw=black!70},
    dashed_arrow/.style={thick, ->, draw=black!50, dashed},
    
    % 背景框样式
    stagebox/.style={rectangle, draw=gray!60, dashed, thick, rounded corners=8pt, inner sep=18pt, fill=gray!2, fill opacity=0.5}
]

% -------------------------
% 1. 节点排列与定义
% -------------------------

% 起点节点
\node (start) [process, text width=5cm] {\textbf{LLM 完整响应}\\ \textit{(CoT 推理过程 + 最终表达式)}};

% 阶段一节点
\node (stage1_check) [decision, below=1cm of start, text width=3cm] {是否包含多标签\\或中英文废话?};
\node (stage1_fail) [penalty, right=1.5cm of stage1_check] {
    \centerline{\textbf{\textcolor{red!80!black}{严重违规/直接截断}}}\\[1mm]
    \textit{例:} "输出为：$<$think$>$"\\[0.5mm]
    \textbf{得分:} $\mathbf{R = -1.0}$
};

% 阶段二节点
\node (stage2_check) [decision, below=2cm of stage1_check, text width=3cm] {思维链推演($A\!=\!B$)\\是否存在幻觉捏造?};
\node (stage2_fail) [penalty, right=2.0cm of stage2_check] {
    \centerline{\textbf{\textcolor{red!80!black}{过程态作弊/查杀}}}\\[1mm]
    \textit{例:} "先算 $3 \times 6 = 12$"\\[0.5mm]
    \textbf{得分:} $\mathbf{R_{\text{step}} = -1.0}$
};
\node (stage2_pass) [reward, left=1.5cm of stage2_check] {
    \centerline{\textbf{\textcolor{green!40!black}{逻辑自洽/密度牵引}}}\\[1mm]
    \textit{例:} "$3 \times 6 = 18$"\;(\textcolor{green!60!black}{\checkmark})\\[0.5mm]
    \textbf{得分:} $\mathbf{+0.1}$ / 每合法步
};

% 提取结果汇总节点
\node (hub) [process, below=2cm of stage2_check] {\textbf{提取待测运算式}\\ \textit{(抛弃前文，严格解析尾部字符串)}};

% 阶段三节点
\node (stage3_check) [decision, below=1.2cm of hub, text width=3cm] {多重集余弦校验\\与无理异常监控};
\node (stage3_invalid) [penalty, left=1.5cm of stage3_check] {
    \centerline{\textbf{\textcolor{red!80!black}{无效穷举/破坏约束}}}\\[1mm]
    \textit{例:} 多用了目标外数字\\[0.5mm]
    \textbf{得分:} $\mathbf{R \in [-0.6,\, -0.2]}$
};
\node (stage3_eval) [decision, below=2cm of stage3_check, text width=2.5cm] {评估值距目标\\偏离度 ($dist$)};
\node (stage3_dense) [dense, left=1.5cm of stage3_eval] {
    \centerline{\textbf{\textcolor{orange!80!black}{连续测度平滑衰减}}}\\[1mm]
    \textit{例:} 算得 22，偏离度为 2\\[0.5mm]
    \textbf{得分:} $0.8 \cdot e^{-0.2} \approx \boldsymbol{+0.65}$
};
\node (stage3_success) [reward, right=1.5cm of stage3_eval] {
    \centerline{\textbf{\textcolor{green!40!black}{精确等价/全局最优}}}\\[1mm]
    \textit{例:} $3 \times 6 + 8 - 2 = 24$\\[0.5mm]
    \textbf{得分:} $\mathbf{R_{\text{final}} = +4.0}$
};


% -------------------------
% 2. 绘制箭头与边
% -------------------------

% 第一阶段连线
\draw [arrow] (start) -- (stage1_check);
\draw [arrow] (stage1_check) -- node[above, font=\scriptsize] {是 (格式崩溃)} (stage1_fail);
\draw [arrow] (stage1_check) -- node[right, font=\scriptsize] {否 (格式合规)} (stage2_check);

% 第二阶段连线
\draw [arrow] (stage2_check) -- node[above, font=\scriptsize] {发生捏造 ($A \neq B$)} (stage2_fail);
\draw [arrow] (stage2_check) -- node[above, font=\scriptsize] {自洽推演} (stage2_pass);
\draw [arrow] (stage2_fail) |- (hub);
\draw [arrow] (stage2_pass) |- (hub);

% 第三阶段连线
\draw [arrow] (hub) -- (stage3_check);
\draw [arrow] (stage3_check) -- node[above, font=\scriptsize] {校验不通过} (stage3_invalid);
\draw [arrow] (stage3_check) -- node[right, font=\scriptsize] {准入白名单} (stage3_eval);
\draw [arrow] (stage3_eval) -- node[above, font=\scriptsize] {$dist \neq 0$} (stage3_dense);
\draw [arrow] (stage3_eval) -- node[above, font=\scriptsize] {$dist = 0$} (stage3_success);

% 最终汇总线 (虚线示意总奖励输出)
\node (final_reward) [process, below=1.8cm of stage3_eval, fill=blue!15, text width=5cm] {\textbf{马尔可夫决策归总奖励}\\ $R_{\text{total}} = R_{\text{step}} + R_{\text{final}}$};
\draw [dashed_arrow] (stage3_invalid.west) -- ++(-0.6,0) |- (final_reward.west);
\draw [dashed_arrow] (stage3_dense.south) |- (final_reward.west);
\draw [dashed_arrow] (stage3_success.south) |- (final_reward.east);


% -------------------------
% 3. 绘制层级背景框
% -------------------------
\begin{scope}[on background layer]
    % 涵盖阶段1
    \node [stagebox, fit=(stage1_check) (stage1_fail)] (box1) {};
    \node [anchor=south west, text=gray!80!black, font=\bfseries, shift={(0.5, 0.1)}] at (box1.north west) {阶段一：动作掩码与早期过滤};

    % 涵盖阶段2
    \node [stagebox, fit=(stage2_check) (stage2_fail) (stage2_pass)] (box2) {};
    \node [anchor=south west, text=gray!80!black, font=\bfseries, shift={(0.5, 0.1)}] at (box2.north west) {阶段二：弱过程监督 (Weak PRM) 测试};

    % 涵盖阶段3
    \node [stagebox, fit=(stage3_check) (stage3_invalid) (stage3_eval) (stage3_dense) (stage3_success)] (box3) {};
    \node [anchor=south west, text=gray!80!black, font=\bfseries, shift={(0.5, 0.1)}] at (box3.north west) {阶段三：终态多重集验证与连续距离反馈};
\end{scope}

\end{tikzpicture}
\end{document}